% ====================================================================================
\chapter*{附录：核心概念索引}
\addcontentsline{toc}{chapter}{附录：核心概念索引}
% ====================================================================================

本附录提供全书核心概念的快速索引，帮助读者：
\begin{itemize}
    \item 查找概念的首次出现位置和完整定义
    \item 理解同一概念在不同章节中的不同形态
    \item 复习时快速定位关键内容
\end{itemize}

% ------------------------------------------------------------------------------------
\section*{核心概念跨章对照表}

\begin{center}
\renewcommand{\arraystretch}{1.6}
\small
\begin{tabular}{p{2.8cm}|p{2.2cm}|p{8cm}}
\toprule
\textbf{概念} & \textbf{首次出现} & \textbf{各章节中的形态} \\
\midrule

\textbf{拉格朗日乘子} $\lambda$
& 第1章
& \textbf{Ch1}: 约束的"影子价格" \newline
  \textbf{Ch5}: 对偶问题的变量 \newline
  \textbf{Ch6}: 最大熵分布的特征权重 \newline
  \textbf{Ch8}: 伴随变量（协态变量） \newline
  \textbf{Ch9}: 反向传播中的误差信号 \newline
  \textbf{Ch15}: 博弈论中的边际价值 \\
\midrule

\textbf{KKT条件}
& 第1章
& \textbf{Ch1}: 约束优化的必要条件 \newline
  \textbf{Ch5}: 强对偶性的充分条件 \newline
  \textbf{Ch7}: SVM的支持向量条件 \newline
  \textbf{Ch15}: 纳什均衡的刻画 \\
\midrule

\textbf{对偶性}
& 第5章
& \textbf{Ch5}: 原问题与对偶问题的关系 \newline
  \textbf{Ch6}: 最大熵 = 最小KL散度 \newline
  \textbf{Ch7}: SVM原问题与对偶问题 \\
\midrule

\textbf{梯度/导数}
& 第3章
& \textbf{Ch3}: 矩阵求导规则 \newline
  \textbf{Ch4}: 正则化项的梯度 \newline
  \textbf{Ch9}: 反向传播计算梯度 \newline
  \textbf{Ch16}: 可微物理中的梯度 \\
\midrule

\textbf{伴随方程/变量}
& 第8章
& \textbf{Ch8}: 最优控制中的协态方程 \newline
  \textbf{Ch9}: 神经网络的反向传播 \newline
  \textbf{Ch16}: 可微仿真中的梯度计算 \\
\midrule

\textbf{熵}
& 第6章
& \textbf{Ch6}: Shannon熵，信息量的期望 \newline
  \textbf{Ch6}: KL散度中的负熵项 \newline
  \textbf{Ch12-5}: 注意力中的熵惩罚 \\
\midrule

\textbf{Softmax}
& 第6章
& \textbf{Ch6}: 最大熵分布的形式 \newline
  \textbf{Ch9}: 分类器输出层 \newline
  \textbf{Ch12}: Attention的权重归一化 \\
\midrule

\textbf{注意力机制}
& 第12章
& \textbf{Ch12}: Query-Key-Value结构 \newline
  \textbf{Ch12-5}: FlashAttention优化 \newline
  \textbf{Ch12-5}: 多头注意力 \\
\midrule

\textbf{正则化}
& 第4章
& \textbf{Ch4}: L1/L2范数惩罚 \newline
  \textbf{Ch7}: SVM的间隔最大化 \newline
  \textbf{Ch12-5}: MoE的负载均衡损失 \\
\bottomrule
\end{tabular}
\end{center}

% ------------------------------------------------------------------------------------
\section*{核心概念快速定义}

\subsection*{A. 优化基础（第1-5章）}

\begin{description}
    \item[拉格朗日乘子 (Lagrange Multiplier)]
    约束优化中引入的辅助变量，表示约束放松一个单位时目标函数的改进量（"影子价格"）。
    \[
    \mathcal{L}(x, \lambda) = f(x) + \lambda g(x)
    \]

    \item[KKT条件 (Karush-Kuhn-Tucker Conditions)]
    约束优化问题最优解的必要条件，包括：
    \begin{enumerate}
        \item 梯度条件：$\nabla f + \sum \lambda_i \nabla g_i = 0$
        \item 可行性：$g_i(x) \leq 0$（或 $= 0$）
        \item 互补松弛：$\lambda_i g_i(x) = 0$
        \item 对偶可行：$\lambda_i \geq 0$（不等式约束）
    \end{enumerate}

    \item[对偶问题 (Dual Problem)]
    通过拉格朗日函数构造的另一个优化问题，提供原问题最优值的下界。强对偶性成立时，两者最优值相等。

    \item[凸优化 (Convex Optimization)]
    目标函数和可行域都是凸的优化问题。局部最优 = 全局最优，KKT条件是充分必要条件。
\end{description}

\subsection*{B. 信息与学习（第6-7章）}

\begin{description}
    \item[熵 (Entropy)]
    分布不确定性的度量，也是平均信息量：
    \[
    H(p) = -\sum_i p_i \log p_i
    \]

    \item[KL散度 (Kullback-Leibler Divergence)]
    两个分布之间的"差异"（非对称）：
    \[
    D_{KL}(P \| Q) = \sum_i p_i \log \frac{p_i}{q_i}
    \]

    \item[交叉熵 (Cross Entropy)]
    用分布 $Q$ 编码分布 $P$ 的平均比特数：
    \[
    H(P, Q) = -\sum_i p_i \log q_i = H(P) + D_{KL}(P \| Q)
    \]

    \item[Softmax函数]
    将实数向量转换为概率分布的函数，是最大熵分布的形式：
    \[
    p_i = \frac{e^{z_i}}{\sum_j e^{z_j}}
    \]

    \item[支持向量机 (SVM)]
    间隔最大化的分类器，对偶问题只依赖于支持向量。
\end{description}

\subsection*{C. 控制与动态（第8章、第16章）}

\begin{description}
    \item[最优控制 (Optimal Control)]
    选择控制输入 $u(t)$ 使得系统轨迹最优化某个目标泛函。

    \item[伴随方程 (Adjoint Equation)]
    描述"未来代价对当前状态的敏感度"的微分方程，从终点向起点反向求解：
    \[
    \dot{\lambda} = -\frac{\partial H}{\partial x}
    \]

    \item[哈密顿函数 (Hamiltonian)]
    综合目标函数和约束的函数：
    \[
    H(x, u, \lambda) = L(x, u) + \lambda^\top f(x, u)
    \]

    \item[可微仿真 (Differentiable Simulation)]
    物理仿真的可微分版本，允许通过反向传播计算控制输入的梯度。
\end{description}

\subsection*{D. 神经网络（第9章、第12章）}

\begin{description}
    \item[反向传播 (Backpropagation)]
    高效计算神经网络梯度的算法，本质是伴随方程的离散版本。

    \item[计算图 (Computational Graph)]
    表示计算过程的有向无环图，前向传播计算值，反向传播计算梯度。

    \item[Transformer]
    基于自注意力机制的神经网络架构，不使用循环或卷积。

    \item[注意力机制 (Attention)]
    根据Query和Key的相似度对Value加权求和：
    \[
    \text{Attention}(Q, K, V) = \text{softmax}\left(\frac{QK^\top}{\sqrt{d_k}}\right) V
    \]

    \item[FlashAttention]
    IO感知的注意力算法，通过分块计算减少内存访问。

    \item[MoE (Mixture of Experts)]
    稀疏激活的架构，每个输入只激活部分专家网络。
\end{description}

\subsection*{E. 多智能体（第15章）}

\begin{description}
    \item[纳什均衡 (Nash Equilibrium)]
    每个玩家都采用最优响应的策略组合，没有人有动机单方面偏离。

    \item[最优响应 (Best Response)]
    给定其他玩家策略时，最大化自己收益的策略。

    \item[Cournot竞争]
    企业同时选择产量的寡头竞争模型，纳什均衡通过KKT条件求解。
\end{description}

% ------------------------------------------------------------------------------------
\section*{概念演化图谱}

\begin{center}
\begin{tikzpicture}[
    node distance=1.8cm,
    concept/.style={rectangle, rounded corners, draw=blue!50, fill=blue!10,
                    text width=2.5cm, align=center, minimum height=1cm},
    arrow/.style={->, >=stealth, thick, blue!60}
]

% 第一行：基础
\node[concept] (lagrange) {拉格朗日乘子\\(Ch1)};
\node[concept, right=of lagrange] (kkt) {KKT条件\\(Ch1)};
\node[concept, right=of kkt] (dual) {对偶性\\(Ch5)};

% 第二行：分支应用
\node[concept, below=of lagrange] (adjoint) {伴随方程\\(Ch8)};
\node[concept, below=of kkt] (entropy) {最大熵\\(Ch6)};
\node[concept, below=of dual] (svm) {SVM\\(Ch7)};

% 第三行：现代应用
\node[concept, below=of adjoint] (bp) {反向传播\\(Ch9)};
\node[concept, below=of entropy] (softmax) {Softmax\\(Ch6,12)};
\node[concept, below=of svm] (nash) {纳什均衡\\(Ch15)};

% 第四行：前沿
\node[concept, below=of bp] (diffsim) {可微仿真\\(Ch16)};
\node[concept, below=of softmax] (attention) {注意力\\(Ch12)};

% 箭头
\draw[arrow] (lagrange) -- (kkt);
\draw[arrow] (kkt) -- (dual);
\draw[arrow] (lagrange) -- (adjoint);
\draw[arrow] (lagrange) -- (entropy);
\draw[arrow] (dual) -- (svm);
\draw[arrow] (adjoint) -- (bp);
\draw[arrow] (entropy) -- (softmax);
\draw[arrow] (kkt) -- (nash);
\draw[arrow] (bp) -- (diffsim);
\draw[arrow] (softmax) -- (attention);

% 虚线连接（表示概念联系）
\draw[dashed, gray] (adjoint) -- (bp);
\draw[dashed, gray] (bp) to[bend right=20] (diffsim);

\end{tikzpicture}
\end{center}

\vspace{0.5cm}
\begin{center}
\textit{箭头表示概念的继承/发展关系，虚线表示本质相同的不同表述。}
\end{center}

% ------------------------------------------------------------------------------------
\section*{本书的核心信息}

\begin{unification}[title=一个公式，多个化身]
拉格朗日函数的梯度条件
\[
\nabla_x \mathcal{L} = \nabla f + \lambda \nabla g = 0
\]
在不同领域有不同的名字：

\begin{center}
\renewcommand{\arraystretch}{1.4}
\begin{tabular}{l|l}
\toprule
\textbf{领域} & \textbf{名字} \\
\midrule
微积分 & 拉格朗日乘子法 \\
运筹学 & KKT条件 \\
物理学 & 哈密顿原理 \\
控制论 & 伴随方程 \\
机器学习 & 反向传播 \\
信息论 & 最大熵分布 \\
博弈论 & 纳什均衡条件 \\
\bottomrule
\end{tabular}
\end{center}

\textbf{理解了拉格朗日乘子的本质，就等于理解了它们的共同根源。}
\end{unification}

% ------------------------------------------------------------------------------------
\section*{章节速查}

\begin{center}
\renewcommand{\arraystretch}{1.3}
\begin{tabular}{c|l|l}
\toprule
\textbf{章} & \textbf{标题} & \textbf{核心关键词} \\
\midrule
1 & 从高数到约束优化 & 拉格朗日乘子、KKT条件、影子价格 \\
2 & 线性代数回顾 & 特征值、SVD、正定矩阵 \\
3 & 矩阵求导 & 链式法则、Jacobian、Hessian \\
4 & 范数与正则化 & L1/L2范数、稀疏性、过拟合 \\
5 & 凸优化与对偶 & 对偶问题、强对偶、凸集 \\
6 & 信息与编码 & 熵、KL散度、交叉熵、Softmax \\
7 & 支持向量机 & 间隔、核技巧、对偶形式 \\
8 & 最优控制 & 伴随方程、Pontryagin原理 \\
9 & 神经网络 & 反向传播、计算图、自动微分 \\
10 & 动态规划 & Bellman方程、值函数 \\
12 & Transformer & 自注意力、多头注意力 \\
12-5 & 高效计算 & FlashAttention、MoE \\
14 & 强化学习 & 策略梯度、Actor-Critic \\
15 & 多智能体 & 纳什均衡、博弈论 \\
16 & 世界模型 & 可微仿真、模型预测控制 \\
\bottomrule
\end{tabular}
\end{center}
